\chapter{Introduction and Motivation}

\comment{introduzione deve essere scritta alla fine}

\section{Project context (taken from \texttt{Methods\_GC\_LD.pdf})}
\comment{Thesis proposal, ecoinformatics, Lucia.}

Eco-informatics is a modern way of conducting scientific research that enables scientists to generate
new knowledge through innovative tools and approaches. Ecology is increasingly becoming a data-intensive
science that involves acquiring, storing, and analysing large amounts of data from a range of sources,
including satellite and aerial remote sensing, instruments, sensors, and human observations. Ecologists
increasingly address large-scale questions with both scientific and social relevance, such as climate change
and its future effects on ecosystems and the ecosystem services associated with them. The topic of this
thesis project originated from a research study conducted by the Department of Earth, Environmental and Life
Sciences (DISTAV, University of Genoa). The study analyses changes in the distribution of natural habitats
in the Alps caused by future climate change in order to support decision-making on conservation planning and
prioritisation. The analysis uses Species Distribution Models, a dataset of approximately four million
records, and future climate-change scenarios based on General Circulation Models. The statistical software
R (https://www.r-project.org/) is used to generate predictions of future changes in habitat distribution and
biodiversity loss. The thesis student will improve the working script and its algorithmic functions to make
the model analysis more efficient, drawing on their knowledge and studies in computer science. This thesis
project involves interdisciplinary collaboration between DIBRIS and DISTAV, which can help produce better
scientific results.
